\begin{center}
    \LARGE \textbf{Хичээлийн танилцуулга} \\
    \vspace{0.3cm}
    \Large \textbf{Хичээлийн нэр: Математик 1B (Advanced Calculus I)} \\
    \large Түвшин: Бакалаврын 1-р курс
\end{center}

\section*{Зорилго}
Оюутнуудад математикийн суурь ойлголтууд болох олонлог, функц, функцийн хязгаар, уламжлал болон интегралын тухай үндсэн ойлголтыг өгөх, тэдгээрийг бодит амьдрал болон техникийн асуудлыг шийдвэрлэх, судалгаа шинжилгээний суурьт ашиглаж сургахад оршино.
\section*{Хичээлийн төлөвлөгөө (Syllabus)}
\renewcommand{\arraystretch}{1.5} % Мөрийн хоорондох зайг нэмэх
\begin{table}[ht]
\centering
\small
\begin{tabular}{|c|p{4cm}|p{8cm}|}
\hline
\textbf{Долоо хоног} & \textbf{Сэдэв} & \textbf{Гол ойлголтууд} \\ \hline
1-2 & Олонлог, функц & Олонлог, баталгааны аргууд, функц, үндсэн функцүүд, график, чанарууд. \\ \hline
3-4 & Хязгаар ба Тасралтгүй чанар & Хязгаарын тодорхойлолт, хязгаар бодох техникүүд, функцийн тасралтгүй чанар. \\ \hline
5-6 & Уламжлалын үндэс & Уламжлалын тодорхойлолт, геометрийн утга (шүргэгч шулуун), дифференциалчлах дүрмүүд. \\ \hline
7-8 & Уламжлал авах техникүүд & Давхар функцийн уламжлал (Chain Rule), далд функцийн уламжлал, тригонометр функцийн уламжлал. \\ \hline
9 & \textbf{Явцын шалгалт (Midterm)} & Өмнөх сэдвүүдийн бататгал. \\ \hline
10-11 & Уламжлалын хэрэглээ & Функцийн өсөх, буурах муж, экстремум цэгүүд (Maximum/Minimum), Лопиталын дүрэм. \\ \hline
12-13 & Интеграл тоолол & Эх функц, тодорхой бус интеграл, орлуулах арга. \\ \hline
14-15 & Тодорхой интеграл & Ньютон-Лейбницийн томьёо, талбай болон эзлэхүүн олох. \\ \hline
16 & Интегралын хэрэглээ & Нумын урт, эргэлтийн биеийн эзлэхүүн, физик хэрэглээ (ажил, хүч). \\ \hline
\end{tabular}
\end{table}
\section*{Ашиглах материалууд}
\begin{itemize}[label=--]
    \item \textbf{Үндсэн сурах бичиг:} James Stewart, "Calculus: Early Transcendentals"
    \item \textbf{Нэмэлт эх сурвалж:} Khan Academy, MIT OpenCourseWare
    \item \textbf{Програм хангамж:} Desmos, GeoGebra, WolframAlpha
\end{itemize}

\newpage
\section*{Үнэлгээний журам (108 оноо)}

\subsection*{Ирц — 16 оноо}
\begin{itemize}
\item Лекцийн ирц: 10 оноо (1 тасалбал 1 оноо хасна)
\item Семинарын ирц: 6 оноо (1 тасалбал 1 оноо хасна)
\item Лекц эсвэл семинарыг \textbf{6 удаа тасалбал} \textbf{WF} үнэлгээ авна
\end{itemize}

\subsection*{Семинарын оноо  — 20 оноо}
Хичээлийн явцад семинарын хичээл дээр 5 удаа самбарт бодлого бодож  аман асуултад хариулна.\\
Нэг удаагийн асуулт: 5 хүртэл оноо\\ Асуултын агуулга:
\begin{itemize}
\item Олонлог, баталгааны аргууд, функц, үндсэн
функцүүд, график, чанарууд, хязгаар бодох
техникүүд, функцийн тасралтгүй чанар.
\item Уламжлалын тодорхойлолт, геометрийн утга
(шүргэгч шулуун), дифференциалчлах дүрмүүд, давхар функцийн уламжлал (Chain Rule),  тригонометр функцийн уламжлал.
\item Функцийн өсөх, буурах муж, экстремум цэгүүд
(Maximum/Minimum), Лопиталын дүрэм, Функцийн шинжилгээ
\item Эх функц, тодорхой бус интеграл, орлуулах арга.
\item Ньютон-Лейбницийн томьёо, талбай болон
эзлэхүүн олох.
\end{itemize}

\subsection*{Гэрийн даалгавар — 7 оноо}

\subsection*{Явцын шалгалт — 25 оноо}
\begin{itemize}
\item 9-р долоо хоногт авна
\item Эхний 8 лекцийн агуулга 
\end{itemize}

\subsection*{Бие даалт — 10 оноо}
\begin{itemize}
\item 12-р долоо хоногт

\end{itemize}

\subsection*{Улиралын шалгалт — 30 оноо}
\begin{itemize}
\item 17-р долоо хоногт авна
\item Сүүлийн 9 лекцийн агуулга (интегралаас хойших)
\end{itemize}

\subsection*{Нийт оноо}
108 онооноос авсан оноогоор дүн гаргана.  
100-аас дээш оноо авсан тохиолдолд 100 гэж дүгнэнэ.